% !TeX root = ../thesis.tex
\chapter{原型系统设计与实现}

\section{系统总体设计}
虽然本文构建的知识抽取模型库尚未最终并入正式生产平台，但已经形成可独立运行的原型系统，用于承接模型登记、制品管理、部署创建、Profile 配置、基准测试和运行观测等关键环节。与单纯展示若干训练脚本或实验结果不同，本章更关注模型库在系统层面的组织方式，即如何把不同知识抽取模型转化为结构统一、接口统一、可部署、可回归和可观察的软件资产。

从工程目标看，该原型系统主要服务于三类需求。第一，面向模型库构建，需要对 CasRel、SpERT、TPLinker、PURE、PRGC、OneRel、CopyMTL、DyGie 和 GraphRel 等模型进行统一登记，保存模型名称、版本、语言、主干编码器、标签映射和 ONNX 制品路径等元数据。第二，面向部署验证，需要将模型制品接入 GPUStack 的本地部署流程，完成实例创建、后端绑定、节点调度和健康检查。第三，面向运维与选型，需要通过 Profile 驱动的基准测试和运行监控，对不同模型在延迟、吞吐量、显存占用和稳定性上的表现进行持续比较。

\begin{table}[htbp]
  \centering
  \caption{原型系统的主要设计目标}
  \vspace{0.5em}\wuhao
  \begin{tabularx}{\textwidth}{l X X}
    \toprule
    设计目标 & 目标说明 & 对论文工作的支撑作用 \\
    \midrule
    模型统一管理 & 统一保存模型制品、元数据、标签映射、语言信息和版本记录 & 支撑知识抽取模型库构建与持续扩充 \\
    工程化部署 & 支持本地模型上传、实例创建、资源绑定和健康检查 & 支撑模型从实验环境向平台原型迁移 \\
    配置化优化 & 通过 Profile 模板管理 batch、并发、精度和线程参数 & 支撑部署优化与横向性能比较 \\
    运行可观测 & 输出日志、指标、任务历史和导出记录 & 支撑问题定位、回溯分析和回滚决策 \\
    \bottomrule
  \end{tabularx}
  \label{tab:system_targets_compact}
\end{table}

为降低不同模型接入时的异构性，本文在系统层采用统一制品描述方式。设某一模型版本的系统级描述对象为
\begin{equation}
\mathcal{A}=\langle id,\ name,\ lang,\ backbone,\ schema,\ artifact,\ profile \rangle ,
\label{eq:artifact_object_system}
\end{equation}
其中 $id$ 表示模型版本唯一标识，$name$ 表示模型名称，$lang$ 表示语言类型，$backbone$ 表示主干编码器，$schema$ 表示标签与输入输出约束，$artifact$ 表示 ONNX 制品及相关文件集合，$profile$ 表示默认部署配置。该定义保证了不同模型虽然内部结构各异，但在平台侧能够以统一对象被读取、展示与部署。

在系统路由层，本文仍采用简化的场景匹配策略。设场景描述为 $\mathbf{s}$，模型为 $m$，则推荐分数可定义为
\begin{equation}
\Psi(m,\mathbf{s})=\omega_{1}F_{1}(m)+\omega_{2}G(m,\mathbf{s})-\omega_{3}L_{95}(m),
\label{eq:scene_route_compact}
\end{equation}
其中 $F_{1}(m)$ 表示模型离线效果，$G(m,\mathbf{s})$ 表示模型与场景的匹配度，$L_{95}(m)$ 表示尾延迟指标。该式并不用于构造复杂调度算法，而是说明系统在展示“推荐部署模型”时遵循的可解释逻辑。

\begin{figure}[htbp]
  \centering
  \begin{minipage}[t]{0.48\textwidth}
    \centering
    \includegraphics[width=\textwidth]{gpustack_screens/01_overview.png}
    \vspace{0.3em}

    \small （a）系统总览页面
  \end{minipage}\hfill
  \begin{minipage}[t]{0.48\textwidth}
    \centering
    \includegraphics[width=\textwidth]{gpustack_screens/19_task_history.png}
    \vspace{0.3em}

    \small （b）任务历史页面
  \end{minipage}
  \caption{原型系统总览与全局任务入口}
  \label{fig:system_overview_compact}
\end{figure}

\section{模型库管理模块实现}
模型库管理模块是本研究与论文主题联系最紧密的系统部分，其任务是把不同知识抽取模型由“训练结果”转化为“可管理制品”。在原型系统中，该模块主要包含模型列表、模型详情、本地上传和模型文件浏览四类页面，用于承接模型注册、版本维护和制品整理等工作。对于企业知识平台而言，这一模块的意义不只是展示模型名称，而是保证平台能够明确知道每个模型的语言版本、依赖编码器、输入输出模式以及对应的部署约束。

由于本研究同时涉及中文模型与英文模型，因此制品组织方式必须保持双语一致。具体而言，中文版本统一采用 bert-base-chinese 作为主干编码器；英文版本则保留英文主干编码器配置，如 bert-base-uncased 或原始英文实现所依赖的 BERT 编码器。但无论是中文还是英文，二者在目录结构、元数据字段、ONNX 导出结果、标签文件命名和服务封装接口上均采用同一套格式，从而实现“语言不同、制品描述一致”的管理效果。

\begin{table}[htbp]
  \centering
  \caption{中英文模型统一制品字段示例}
  \vspace{0.5em}\wuhao
  \begin{tabularx}{\textwidth}{l l X}
    \toprule
    字段名称 & 示例取值 & 说明 \\
    \midrule
    model\_id & casrel-zh-v1 / casrel-en-v1 & 区分语言版本和模型版本 \\
    language & zh / en & 标识中文或英文模型 \\
    backbone & bert-base-chinese / bert-base-uncased & 显式记录对应主干编码器 \\
    artifact\_path & models/casrel/\{lang\}/v1/model.onnx & 中英文共享同一目录组织规则 \\
    schema\_path & meta/schema.json & 统一保存标签与输出结构定义 \\
    tokenizer\_path & meta/tokenizer/ & 与语言版本绑定，但字段格式一致 \\
    default\_profile & profiles/default.yaml & 统一引用默认部署模板 \\
    \bottomrule
  \end{tabularx}
  \label{tab:bilingual_artifact_schema}
\end{table}

\begin{figure}[htbp]
  \centering
  \begin{minipage}[t]{0.48\textwidth}
    \centering
    \includegraphics[width=\textwidth]{gpustack_screens/02_models.png}
    \vspace{0.3em}

    \small （a）模型库列表页面
  \end{minipage}\hfill
  \begin{minipage}[t]{0.48\textwidth}
    \centering
    \includegraphics[width=\textwidth]{gpustack_screens/03_model_detail.png}
    \vspace{0.3em}

    \small （b）模型详情页面
  \end{minipage}
  \caption{模型库管理模块界面一}
  \label{fig:model_module_ui_1}
\end{figure}

\begin{figure}[htbp]
  \centering
  \begin{minipage}[t]{0.48\textwidth}
    \centering
    \includegraphics[width=\textwidth]{gpustack_screens/04_upload_model.png}
    \vspace{0.3em}

    \small （a）本地模型上传页面
  \end{minipage}\hfill
  \begin{minipage}[t]{0.48\textwidth}
    \centering
    \includegraphics[width=\textwidth]{gpustack_screens/13_model_files.png}
    \vspace{0.3em}

    \small （b）模型文件浏览页面
  \end{minipage}
  \caption{模型库管理模块界面二}
  \label{fig:model_module_ui_2}
\end{figure}

从系统实现角度看，模型库管理模块完成了三项关键工作。其一，统一模型元数据结构，使中文与英文模型在平台层具有相同的注册方式。其二，统一本地模型上传入口，使导出的 ONNX 制品能够被平台直接识别，而不再依赖分散的实验目录。其三，统一模型文件与版本浏览方式，使后续部署、回滚和排错都可以追溯到具体模型版本。这些工作共同构成了模型工程化部署的入口层。

\section{部署编排模块实现}
部署编排与资源调度模块的任务，是将模型库中的静态制品转换为可调用的推理服务实例。该模块包含部署列表、部署创建、部署详情、服务路由、资源提供方、计算节点、GPU 状态和推理后端等页面，覆盖了“模型选择、实例启动、资源绑定、运行检查”的完整链路。对于本研究而言，这一模块比单纯的模型精度对比更重要，因为它直接决定模型是否具备平台接入能力。

在部署创建过程中，系统要求同时明确模型版本、运行节点、推理后端、精度模式、资源额度和健康检查策略。若记模型实例 $m$ 在节点 $n$ 上的部署方案为 $d=(m,n,p,b)$，其中 $p$ 为 Profile 配置，$b$ 为后端类型，则部署可行性可抽象表示为
\begin{equation}
\Omega(d)=\mathbb{I}\!\left(\mathrm{Mem}_{req}(m,p)\leq \mathrm{Mem}_{free}(n)\right)
\cdot
\mathbb{I}\!\left(\mathrm{Compat}(m,b)=1\right),
\label{eq:deploy_feasibility_view}
\end{equation}
其中 $\mathrm{Mem}_{req}(m,p)$ 表示模型在给定 Profile 下的显存需求，$\mathrm{Mem}_{free}(n)$ 表示节点可用显存，$\mathrm{Compat}(m,b)$ 表示模型与推理后端之间的兼容关系。这表明部署界面中看似简单的参数输入，实际对应资源、后端和模型三者之间的联合约束。

\begin{figure}[htbp]
  \centering
  \begin{minipage}[t]{0.48\textwidth}
    \centering
    \includegraphics[width=\textwidth]{gpustack_screens/05_deployments.png}
    \vspace{0.3em}

    \small （a）部署列表页面
  \end{minipage}\hfill
  \begin{minipage}[t]{0.48\textwidth}
    \centering
    \includegraphics[width=\textwidth]{gpustack_screens/06_deploy_create.png}
    \vspace{0.3em}

    \small （b）部署创建页面
  \end{minipage}
  \caption{部署编排模块界面一}
  \label{fig:deploy_module_ui_1}
\end{figure}

\begin{figure}[htbp]
  \centering
  \begin{minipage}[t]{0.48\textwidth}
    \centering
    \includegraphics[width=\textwidth]{gpustack_screens/07_deploy_detail.png}
    \vspace{0.3em}

    \small （a）部署详情页面
  \end{minipage}\hfill
  \begin{minipage}[t]{0.48\textwidth}
    \centering
    \includegraphics[width=\textwidth]{gpustack_screens/08_routes.png}
    \vspace{0.3em}

    \small （b）服务路由页面
  \end{minipage}
  \caption{部署编排模块界面二}
  \label{fig:deploy_module_ui_2}
\end{figure}

\begin{figure}[htbp]
  \centering
  \begin{minipage}[t]{0.48\textwidth}
    \centering
    \includegraphics[width=\textwidth]{gpustack_screens/09_providers.png}
    \vspace{0.3em}

    \small （a）资源提供方页面
  \end{minipage}\hfill
  \begin{minipage}[t]{0.48\textwidth}
    \centering
    \includegraphics[width=\textwidth]{gpustack_screens/10_nodes.png}
    \vspace{0.3em}

    \small （b）计算节点页面
  \end{minipage}
  \caption{资源接入与节点管理界面}
  \label{fig:deploy_module_ui_3}
\end{figure}

\begin{figure}[htbp]
  \centering
  \begin{minipage}[t]{0.48\textwidth}
    \centering
    \includegraphics[width=\textwidth]{gpustack_screens/11_gpus.png}
    \vspace{0.3em}

    \small （a）GPU 监测页面
  \end{minipage}\hfill
  \begin{minipage}[t]{0.48\textwidth}
    \centering
    \includegraphics[width=\textwidth]{gpustack_screens/12_backends.png}
    \vspace{0.3em}

    \small （b）推理后端页面
  \end{minipage}
  \caption{资源状态与后端配置界面}
  \label{fig:deploy_module_ui_4}
\end{figure}

本文保留 GPUStack 的基础管理能力，但不引入大模型在线生成、远程模型市场等与论文主题关联较弱的功能。这样做的目的，是使系统实现始终围绕知识抽取模型库展开，重点验证 ONNX 制品的本地化部署、节点绑定、后端切换和实例治理，而不是扩展为通用大模型服务平台。

\section{基准测试模块实现}
基准测试与 Profile 配置模块服务于论文第四章的部署优化研究，其核心思想是将性能测试从“临时命令行试跑”转变为“配置化、可重复、可对比”的系统流程。在该模块中，研究者可以为不同模型预先定义 Profile 文件，将 batch size、并发数、线程数、精度模式、输入长度和预热轮次等参数固化下来，再批量触发测试任务，以获得一组可用于横向比较的部署数据。原型系统随文保留了 \texttt{profiles/default.yaml}、\texttt{profiles/latency-l1.yaml}、\texttt{profiles/throughput-t1.yaml}、\texttt{profiles/memory-m1.yaml} 和 \texttt{profiles/regression-r1.yaml} 等示例配置，用于支撑页面中的模板选择与测试任务创建。

如果记某一 Profile 配置为
\begin{equation}
P=\langle bs,\ conc,\ prec,\ thr,\ len,\ warm \rangle ,
\label{eq:profile_vector_system}
\end{equation}
其中 $bs$ 表示批大小，$conc$ 表示并发度，$prec$ 表示精度模式，$thr$ 表示线程配置，$len$ 表示输入序列长度，$warm$ 表示预热轮次，则 Profile 文件的价值在于把原本隐含在脚本参数中的测试环境显式化。对于不同模型而言，即使采用同样的 ONNX Runtime 或 TensorRT 后端，其最优配置也未必一致，因此论文中提出“编写多个 Profile 文件进行部署优化”是有工程依据的。

\begin{figure}[htbp]
  \centering
  \begin{minipage}[t]{0.48\textwidth}
    \centering
    \includegraphics[width=\textwidth]{gpustack_screens/14_benchmarks.png}
    \vspace{0.3em}

    \small （a）基准测试列表页面
  \end{minipage}\hfill
  \begin{minipage}[t]{0.48\textwidth}
    \centering
    \includegraphics[width=\textwidth]{gpustack_screens/15_benchmark_create.png}
    \vspace{0.3em}

    \small （b）基准任务创建页面
  \end{minipage}
  \caption{基准测试与任务配置界面}
  \label{fig:benchmark_module_ui_1}
\end{figure}

\begin{figure}[htbp]
  \centering
  \begin{minipage}[t]{0.48\textwidth}
    \centering
    \includegraphics[width=\textwidth]{gpustack_screens/16_benchmark_result.png}
    \vspace{0.3em}

    \small （a）基准测试结果页面
  \end{minipage}\hfill
  \begin{minipage}[t]{0.48\textwidth}
    \centering
    \includegraphics[width=\textwidth]{gpustack_screens/20_profile_templates.png}
    \vspace{0.3em}

    \small （b）Profile 模板页面
  \end{minipage}
  \caption{Profile 模板与结果归档界面}
  \label{fig:benchmark_module_ui_2}
\end{figure}

为了在系统层面对不同 Profile 进行统一比较，本文将综合性能评价抽象为
\begin{equation}
S_{perf}=\eta_{1}\frac{Q}{Q_{\max}}+\eta_{2}\left(1-\frac{L_{95}}{L_{95}^{\max}}\right)+\eta_{3}\left(1-\frac{M_{peak}}{M_{peak}^{\max}}\right),
\label{eq:profile_score_view}
\end{equation}
其中 $Q$ 表示吞吐量，$L_{95}$ 表示 P95 延迟，$M_{peak}$ 表示显存峰值，且 $\eta_{1}+\eta_{2}+\eta_{3}=1$。该式用于说明系统为何需要保留结果页与 Profile 模板页，而不要求在本章中展开大量具体性能数值。本章强调的是“模块设计与支撑能力”，具体性能结果则可在后续补充真实实验数据时再进行替换。

\section{监控服务模块实现}
监控观测与结果服务模块承担系统闭环的收尾工作，其目的不是重复展示模型效果表格，而是说明当模型完成部署后，系统如何提供运行可见性和输出可交付性。一方面，日志页和指标页记录实例启动、推理执行、显存占用、吞吐变化和异常信息，为服务状态判断提供依据；另一方面，抽取结果页和结果导出页则负责承载模型输出、置信度、证据片段以及报告生成等功能，使平台使用者能够直观地检查结果结构是否符合知识库写入要求。

\begin{figure}[htbp]
  \centering
  \begin{minipage}[t]{0.48\textwidth}
    \centering
    \includegraphics[width=\textwidth]{gpustack_screens/17_logs.png}
    \vspace{0.3em}

    \small （a）运行日志页面
  \end{minipage}\hfill
  \begin{minipage}[t]{0.48\textwidth}
    \centering
    \includegraphics[width=\textwidth]{gpustack_screens/18_metrics.png}
    \vspace{0.3em}

    \small （b）运行指标页面
  \end{minipage}
  \caption{监控观测模块界面}
  \label{fig:observe_module_ui}
\end{figure}

\begin{figure}[htbp]
  \centering
  \begin{minipage}[t]{0.48\textwidth}
    \centering
    \includegraphics[width=\textwidth]{gpustack_screens/21_extract_result.png}
    \vspace{0.3em}

    \small （a）抽取结果验证页面
  \end{minipage}\hfill
  \begin{minipage}[t]{0.48\textwidth}
    \centering
    \includegraphics[width=\textwidth]{gpustack_screens/22_export_center.png}
    \vspace{0.3em}

    \small （b）结果导出页面
  \end{minipage}
  \caption{结果服务模块界面}
  \label{fig:result_module_ui}
\end{figure}

若以服务运行闭环表示系统状态，则可将其概括为
\begin{equation}
\mathcal{C}=f(\text{deploy},\ \text{observe},\ \text{verify},\ \text{export}),
\label{eq:service_closed_loop}
\end{equation}
其中 $\text{deploy}$ 表示实例已部署，$\text{observe}$ 表示运行状态可观测，$\text{verify}$ 表示结果结构可校验，$\text{export}$ 表示输出结果可导出。只有四个环节同时成立，模型才真正具备“可发布、可使用、可回溯”的工程属性。这也是本章不再展开具体业务文本结果表格的原因，因为系统章节更应强调服务闭环与模块职责，而非重复实验章节中的指标比较。

\begin{table}[htbp]
  \centering
  \caption{第五章系统模块与截图对应关系}
  \vspace{0.5em}\wuhao
  \begin{tabularx}{\textwidth}{l l X}
    \toprule
    模块名称 & 对应截图编号 & 说明 \\
    \midrule
    总体结构入口 & 01, 19 & 展示系统总览、统一导航与任务历史入口 \\
    模型库管理模块 & 02, 03, 04, 13 & 展示模型列表、详情、本地上传和制品浏览能力 \\
    部署编排与资源调度模块 & 05--12 & 展示部署创建、路由管理、节点绑定和后端配置能力 \\
    基准测试与 Profile 配置模块 & 14, 15, 16, 20 & 展示基准任务、Profile 模板与结果归档能力 \\
    监控观测与结果服务模块 & 17, 18, 21, 22 & 展示运行观测、结果校验和导出能力 \\
    \bottomrule
  \end{tabularx}
  \label{tab:screen_mapping_compact}
\end{table}

\FloatBarrier

\section{本章小结}
本章从原型系统设计目标、模型库管理、部署编排、Profile 配置以及监控结果服务五个方面，说明了知识抽取模型库在系统层面的模块设计与工程化实现方式。与前文偏重模型训练和部署优化机制的章节不同，本章强调的是模型如何以统一制品、统一接口和统一页面组织为可部署、可配置、可观察的软件资产。通过对 GPUStack 风格原型界面的模块化展示，本文构建的知识抽取模型库已经具备面向企业平台进一步集成的系统基础。
